% TOP_PROBLEM: {"problemId": "problem.legendres-conjecture", "canonicalTitle": "Legendre's Conjecture on Primes Between Consecutive Squares", "releaseRank": 321, "primaryDomain": "number_theory_arithmetic_geometry", "primaryDomainLabel": "Number theory & arithmetic geometry", "displayStatus": "open", "releaseStatus": "open"}
% !TeX program = lualatex
% Definition and sources only; this file is not an LLM solution attempt.
\documentclass[11pt]{article}
\usepackage[margin=25mm]{geometry}
\usepackage{fontspec}
\setmainfont{Latin Modern Roman}
\usepackage{amsmath,amssymb,mathtools}
\usepackage{unicode-math}
\setmathfont{Latin Modern Math}
\usepackage{xurl,hyperref}
\hypersetup{colorlinks=true,urlcolor=blue}
\setcounter{secnumdepth}{0}
\setlength{\parindent}{0pt}
\setlength{\parskip}{5pt}
\setlength{\emergencystretch}{3em}
\begin{document}
\section*{\texorpdfstring{Legendre's Conjecture on Primes Between Consecutive Squares}{Legendre's Conjecture on Primes Between Consecutive Squares}}\label{321}
\subsection{Definitions and mathematical statement}
A prime is an integer $p>1$ whose only positive divisors are $1$ and $p$. Legendre's conjecture is
\[
 \forall n\in{\mathbb{Z}}_{\ge1}\ \exists p\text{ prime},
 \qquad n^2<p<(n+1)^2.
\]
Both inequalities are strict. Equivalently, if $\pi(x)=\#\{p\le x:p\text{ prime}\}$, one asks $\pi((n+1)^2)-\pi(n^2)\ge1$ for every $n\ge1$, since the upper endpoint is composite. The interval width is $2n+1$, but the assertion is a prime in every one of these specific intervals, not merely an average statement about intervals of this size.
\par\smallskip\textit{Formulation and concept references:} \href{https://mathworld.wolfram.com/LegendresConjecture.html}{[S1]}.

\subsection{Short English statement}
Is there always a prime strictly between one positive integer's square and the next integer's square?
\subsection{Sources}
\begin{itemize}
\item[S1]Weisstein, E. W., Legendre's Conjecture, From MathWorld - A Wolfram Resource, updated 16 August 2026.
\url{https://mathworld.wolfram.com/LegendresConjecture.html}.
\textit{Formulation source.}
\end{itemize}
\end{document}
